% GENERATED FILE. Edit canonical lessons, then run npm run generate:books.
% canonical-source-hash: fnv1a64:79fccc5e129a1717
% canonical-lessons: KA-C143-ona, KA-C143-bada, KA-C143-srimanta, KA-C143-prasiddha, KA-C143-suraksita

\chapter{Dry, Poor, Rich, Famous, Safe}
\label{ch:ka-dry-poor-rich-famous-safe}
\hlchaptermodality{143}

\begin{chapteropening}
\textbf{By the end of this chapter:} \emph{I can say dry, poor, rich, famous and safe.}

Complete the last lesson of chapter 143: I can say dry, poor, rich, famous and safe.
\end{chapteropening}

\section[\texorpdfstring{\kn{ಒಣ} (oṇa)}{ಒಣ (oṇa)}]{\kn{ಒಣ} (oṇa) — dry}

\label{lesson:KA-C143-ona}



\begin{quote}
Before the new one: say the Kannada for narrow, then the Kannada for wet.
\end{quote}

\subsection*{You'll want to know: \kn{ಒಣ}}
\textbf{\kn{ಒಣ}} — \emph{oṇa} — \textquotedblleft{}dry\textquotedblright{}.

\begin{cousinweb}
\textbf{\kn{ಒಣ}} (\emph{oṇa}) is native Dravidian, from \emph{oṇagu}, \textquotedblleft{}to dry\textquotedblright{}, and DEDR 601 lists \emph{oṇa}, \textquotedblleft{}dryness\textquotedblright{}, itself.
\end{cousinweb}

\begin{grammarlens}[title={What you've built}]
One more word for this chapter.
\end{grammarlens}

\subsection*{Your turn}
\begin{itemize}
\raggedright
  \item \emph{Say it:} \emph{oṇa}
  \item \emph{Say it:} \emph{oṇa}, once more
  \item \emph{Say it:} say \emph{oṇa}
  \item \emph{Recall:} say the Kannada for narrow, then the Kannada for wet, then say \emph{oṇa} again
\end{itemize}

\begin{tcolorbox}[breakable,colback=teal!4,colframe=teal!35!black,title={Before you move on}]
What does \kn{ಒಣ} mean? (\textquotedblleft{}Dry\textquotedblright{}.) Say it once more.
\end{tcolorbox}
\section[\texorpdfstring{\kn{ಬಡ} (baḍa)}{ಬಡ (baḍa)}]{\kn{ಬಡ} (baḍa) — poor}

\label{lesson:KA-C143-bada}



\begin{quote}
Before the new one: say the Kannada for wet, then the Kannada for dry.
\end{quote}

\subsection*{You'll want to know: \kn{ಬಡ}}
\textbf{\kn{ಬಡ}} — \emph{baḍa} — \textquotedblleft{}poor\textquotedblright{}.

\begin{cousinweb}
\textbf{\kn{ಬಡ}} (\emph{baḍa}) is native Dravidian, one word with Tamil \emph{vaṭi}, Telugu \emph{baḍugu} and Malayalam \emph{vaṭu} (DEDR 5222). DEDR lists \emph{baḍa} itself, \textquotedblleft{}poor, weak\textquotedblright{}: the old \emph{v-} has become \emph{b-}.
\end{cousinweb}

\begin{grammarlens}[title={What you've built}]
One more word for this chapter.
\end{grammarlens}

\subsection*{Your turn}
\begin{itemize}
\raggedright
  \item \emph{Say it:} \emph{baḍa}
  \item \emph{Say it:} \emph{baḍa}, once more
  \item \emph{Say it:} say \emph{baḍa}
  \item \emph{Recall:} say the Kannada for wet, then the Kannada for dry, then say \emph{baḍa} again
\end{itemize}

\begin{tcolorbox}[breakable,colback=teal!4,colframe=teal!35!black,title={Before you move on}]
What does \kn{ಬಡ} mean? (\textquotedblleft{}Poor\textquotedblright{}.) Say it once more.
\end{tcolorbox}
\section[\texorpdfstring{\kn{ಶ್ರೀಮಂತ} (śrīmanta)}{ಶ್ರೀಮಂತ (śrīmanta)}]{\kn{ಶ್ರೀಮಂತ} (śrīmanta) — rich}

\label{lesson:KA-C143-srimanta}



\begin{quote}
Before the new one: say the Kannada for dry, then the Kannada for poor.
\end{quote}

\subsection*{You'll want to know: \kn{ಶ್ರೀಮಂತ}}
\textbf{\kn{ಶ್ರೀಮಂತ}} — \emph{śrīmanta} — \textquotedblleft{}rich\textquotedblright{}.

\begin{cousinweb}
\textbf{\kn{ಶ್ರೀಮಂತ}} (\emph{śrīmanta}) is Sanskrit \emph{śrīmat}, \textquotedblleft{}rich, fortunate\textquotedblright{}, from \emph{śrī}, \textquotedblleft{}wealth, prosperity\textquotedblright{}, in the form \emph{śrīmant-} that Sanskrit uses before an ending.
\end{cousinweb}

\begin{grammarlens}[title={What you've built}]
One more word for this chapter.
\end{grammarlens}

\subsection*{Your turn}
\begin{itemize}
\raggedright
  \item \emph{Say it:} \emph{śrīmanta}
  \item \emph{Say it:} \emph{śrīmanta}, once more
  \item \emph{Say it:} say \emph{śrīmanta}
  \item \emph{Recall:} say the Kannada for dry, then the Kannada for poor, then say \emph{śrīmanta} again
\end{itemize}

\begin{tcolorbox}[breakable,colback=teal!4,colframe=teal!35!black,title={Before you move on}]
What does \kn{ಶ್ರೀಮಂತ} mean? (\textquotedblleft{}Rich\textquotedblright{}.) Say it once more.
\end{tcolorbox}
\section[\texorpdfstring{\kn{ಪ್ರಸಿದ್ಧ} (prasiddha)}{ಪ್ರಸಿದ್ಧ (prasiddha)}]{\kn{ಪ್ರಸಿದ್ಧ} (prasiddha) — famous}

\label{lesson:KA-C143-prasiddha}



\begin{quote}
Before the new one: say the Kannada for poor, then the Kannada for rich.
\end{quote}

\subsection*{You'll want to know: \kn{ಪ್ರಸಿದ್ಧ}}
\textbf{\kn{ಪ್ರಸಿದ್ಧ}} — \emph{prasiddha} — \textquotedblleft{}famous\textquotedblright{}.

\begin{cousinweb}
\textbf{\kn{ಪ್ರಸಿದ್ಧ}} (\emph{prasiddha}) is Sanskrit \emph{prasiddha}, \textquotedblleft{}celebrated\textquotedblright{}, taken over whole as a learned loan, a tatsama (CDIAL 8842).
\end{cousinweb}

\begin{grammarlens}[title={What you've built}]
One more word for this chapter.
\end{grammarlens}

\subsection*{Your turn}
\begin{itemize}
\raggedright
  \item \emph{Say it:} \emph{prasiddha}
  \item \emph{Say it:} \emph{prasiddha}, once more
  \item \emph{Say it:} say \emph{prasiddha}
  \item \emph{Recall:} say the Kannada for poor, then the Kannada for rich, then say \emph{prasiddha} again
\end{itemize}

\begin{tcolorbox}[breakable,colback=teal!4,colframe=teal!35!black,title={Before you move on}]
What does \kn{ಪ್ರಸಿದ್ಧ} mean? (\textquotedblleft{}Famous\textquotedblright{}.) Say it once more.
\end{tcolorbox}
\section[\texorpdfstring{\kn{ಸುರಕ್ಷಿತ} (surakṣita)}{ಸುರಕ್ಷಿತ (surakṣita)}]{\kn{ಸುರಕ್ಷಿತ} (surakṣita) — safe}

\label{lesson:KA-C143-suraksita}



\begin{quote}
Before the new one: say the Kannada for rich, then the Kannada for famous.
\end{quote}

\subsection*{You'll want to know: \kn{ಸುರಕ್ಷಿತ}}
\textbf{\kn{ಸುರಕ್ಷಿತ}} — \emph{surakṣita} — \textquotedblleft{}safe\textquotedblright{}.

\begin{cousinweb}
\textbf{\kn{ಸುರಕ್ಷಿತ}} (\emph{surakṣita}) is Sanskrit \emph{surakṣita}, \textquotedblleft{}well protected\textquotedblright{}, taken over whole as a learned loan, a tatsama (CDIAL 13499).
\end{cousinweb}

\begin{grammarlens}[title={What you've built}]
One more word for this chapter.
\end{grammarlens}

\subsection*{Your turn}
\begin{itemize}
\raggedright
  \item \emph{Say it:} \emph{surakṣita}
  \item \emph{Say it:} \emph{surakṣita}, once more
  \item \emph{Say it:} say \emph{surakṣita}
  \item \emph{Recall:} say the Kannada for rich, then the Kannada for famous, then say \emph{surakṣita} again
\end{itemize}

\begin{tcolorbox}[breakable,colback=teal!4,colframe=teal!35!black,title={Before you move on}]
What does \kn{ಸುರಕ್ಷಿತ} mean? (\textquotedblleft{}Safe\textquotedblright{}.) Say it once more.
\end{tcolorbox}
